\documentclass[11pt]{article}
\usepackage[a4paper,margin=28mm]{geometry}
\usepackage{amsmath,amsthm,amssymb}
\usepackage[hidelinks]{hyperref}
\newtheorem{theorem}{Theorem}
\title{An effective pathwidth lower bound for Maximum Cycle Cover}
\author{Atlas\thanks{The proof in this note was found using ChatGPT.}}
\date{18 November 2026}
\begin{document}
\maketitle
In \textsc{Maximum Cycle Cover}, the input is a finite simple undirected
graph, a nonnegative binary integer $r$, and a path decomposition of
width $p$. The question is whether the vertices can be partitioned into
at least $r$ simple cycles. Let $B$ be the complete input length.
A finite nonnegative function $h$ is effectively $o(p\log p)$ if there
is a computable function $T$ such that $h(p)\le \phi(p)/j$
whenever $j\ge1$ and $p\ge T(j)$, where
$\phi(p)=\max\{1,p\log_2(p+2)\}$.

\hypertarget{maximum-cycle-cover-supplied-pathwidth-w1}{}
\begin{theorem}
Assume $\mathrm{FPT}\ne\mathrm{W}[1]$. There is no deterministic
algorithm for \textsc{Maximum Cycle Cover} in time
$2^{h(p)}B^{\mathcal O(1)}$ for an effectively $o(p\log p)$ exponent $h$,
even on connected bipartite graphs of maximum degree eight.
\end{theorem}

We reduce from Clique by the requested size. We group Boolean variables
to bound the order of a permutation encoded by a cycle packing. We then
reduce the vertex degrees and extend the packing to a spanning cover.


\section{From Clique to permutation hitting}

We use the \(\mathrm{W[1]}\)-hardness of Clique parameterized by the requested
size \(k\)~\cite{df}. For a graph with \(n\ge 2\) vertices, use \(k\) blocks of
\(\lceil\log_2 n\rceil\) Boolean variables to encode the chosen vertices.
Forbid each unused code in each block, and for each pair of blocks forbid every
pair of valid codes representing equal or nonadjacent vertices. Each prohibition
is a clause on the affected blocks, falsified exactly by the forbidden pattern.
The formula is satisfiable exactly when the graph has a \(k\)-clique. It has
\(m=k\lceil\log_2 n\rceil\) variables, and its complete length and construction
time are both bounded by \(F(k)n^c\) for computable \(F\) and a fixed \(c\).
We decide trivial Clique inputs directly.

Take an explicit CNF over \(m\) listed variables. Decide instances with fewer
than \(16\) variables or an empty clause directly, and likewise decide a
formula with no clauses. Let
\[
 b=\lfloor (\log_2 m)/2\rfloor,\qquad
 q=2^b,\qquad s=\lceil m/b\rceil,\qquad a=\max\{s,q\}.
\]
Pad the variable list to \(sb\) and partition it into \(s\) groups of size
\(b\). In row \(i\) of a \(a\times a\) table, the first \(q\) columns encode
assignments to group \(i\), for each \(i\le s\). Require a row selection to
hit the set of these \(q\) cells for each real row. For each clause, require
it to hit the set of all real cells whose group assignment satisfies a
literal of the clause. Thus a mapping selecting one column per row hits all
sets exactly when its real rows encode a satisfying assignment. Dummy
rows have no domain constraint. The set family is explicit and has polynomial
size in the formula length and \(m\). We have \(a=O(m/\log m)\).

Use the deterministic constructive coloring family of Lokshtanov, Marx, and
Saurabh~\cite[Lemma 2.8 and its proof]{lms}. It enumerates maps
\(f\colon[a]\times[a]\to[a+1]\) in time
\(2^{O(a\log\log(\max\{a,4\}))}a^{O(1)}\). For any mapping selecting one cell
per row, some \(f\) gives the selected cells distinct colors and gives each
selected cell a color absent from the other cells of its row. Indeed, the
graph formed by a clique on the selected cells and the stars joining each
selection to its row is the cactus-grid graph of that lemma.

For each \(f\), retain the cells whose colors are unique within their rows, and
replace the column of every retained cell by its color. Apply this operation to
every set and additionally require a hit on the retained cells of each row.
Add one unconstrained row to make a square of size \(K=a+1\).
For any permutation hitting these new sets, the row-color choice in each
original row has a unique preimage, which gives an original selection hitting
all old sets. Conversely, for a selection hitting all old sets, the map
\(f\) guaranteed by the coloring family retains every chosen cell and assigns
them distinct colors. Assign the unused color to the new dummy row. This
gives a permutation hitting all new sets. The number of queries and the time
to enumerate the family are \(2^{o(m)}\) up to a fixed polynomial factor in
the explicit formula length, with a computable sublinear exponent bound
\(O(m\log\log m/\log m)\).




\section{Representing the permutation by a cycle packing}
We extend the cycle packing construction of Cygan et al.~\cite{cc,ccpre}
to allow multiple cells in a row of a test set. From a
$K\times K$ table with sets $S_1,\ldots,S_M$ of cells, we will construct
a simple graph with a supplied path decomposition of width $O(K)$
in which the packing threshold is attained exactly when a permutation
hits every set. We reject a member containing an empty set and accept
one having no sets directly. During construction we use multigraphs
and count two parallel edges as a cycle.


\subsection{Sparse capacity gadget}

Let $X=\{v_1,\ldots,v_d\}$ be disjoint host ports. To implement host capacity one, introduce d-1 private vertices $u_1$,...,$u_{d-1}$, and put two parallel edges between $u_i$ and each of $v_i$,$v_{i+1}$. For d=1 introduce nothing. Call the resulting doubled alternating path a sparse selector. It has exactly d-1 private vertices.

\textbf{Matching property.} For every A subset X there is a matching of size min(|A|,d-1) between A and the private vertices. If A=X, match u\_i to v\_i. Otherwise choose a missing port v\_j, match u\_i to v\_i for i<j and to v\_\{i+1\} for i\(\ge\)j; this matches every port other than v\_j. Restrict this matching to A.

Every matching edge gives a 2-cycle using its two parallel edges. Thus whenever host cycles use at most one port, all d-1 private vertices can be covered by d-1 disjoint 2-cycles on the remaining ports. Conversely the gadget can never contribute more than d-1 cycles, since each such cycle must use a private vertex.

Also use a capacity-(d-1) gadget on X: one private vertex with two parallel edges to each port. Its full contribution is one 2-cycle, and it requires at least one port unused by host cycles. This definition includes d=1, where the host capacity is zero.

\textbf{Normalization lemma.} Suppose the sets of host ports of all gadgets are pairwise disjoint, and gadget-private vertices have no other neighbors. There is a maximum cycle packing in which every gadget-private vertex lies on a 2-cycle and the remaining host-only packing respects each gadget's capacity.

Proof for one sparse selector, preserving already normalized gadgets. Relative to the current gadget, a kept or host cycle means any cycle avoiding its private vertices; it may still touch other unprocessed gadgets. Such cycles are allowed to be deleted below. Previously normalized private 2-cycles cannot be affected because the port sets are disjoint. Keep every cycle not touching its private vertices. Let c be the number of removed cycles and let A be the ports not used by the kept cycles. Each removed cycle contains at least one private vertex and at least one distinct port; hence c\(\le\)min(d-1,|A|). The matching property supplies c or more 2-cycles, so replacing removed cycles does not decrease the packing size. Now all cycles touching this gadget are 2-cycles; write a for the number of ports used by host cycles and q for the number of these gadget 2-cycles. We have q\(\le\)d-a. If a\(\le\)1, augment the matching to d-1 2-cycles, using the matching property on all ports available to this gadget. If a\(\ge\)2, delete host cycles until at most one port remains used. This costs at most a-1 cycles: if every affected host cycle uses one port, keep one; if some cycle uses at least two, deleting all affected cycles costs at most a-1. Now replace the q gadget cycles by d-1 gadget cycles. The gain d-1-q is at least a-1, compensating every deleted host cycle. Capacity one and full gadget contribution are achieved.

For the one-private-vertex gadget, if its private vertex is not on a 2-cycle, replace the cycle containing it by a 2-cycle, or add a 2-cycle if it was unused and some port is free. If all ports are occupied by host cycles, delete one such cycle and use one of its freed ports for the gadget 2-cycle. This loses at most one cycle and gains one. The resulting host packing uses at most d-1 ports. As gadget port sets are disjoint, these operations do not delete any previously normalized gadget 2-cycle. Process the finitely many gadgets sequentially. Starting from a maximum packing, every nondecreasing replacement preserves maximum size, proving the lemma.

\subsection{Host graph}

For each stage $s$, let $E_s$ consist of the triples $e=(i,j,s)$ for which $(i,j)\in S_s$. Distinct names designate distinct vertices.

Create global vertices $P_i$,$Q_i$ for rows i and $A_j$,$B_j$ for columns j. For every pair i,j create selector ports $p_{ij}$,$q_{ij}$, with edges

\begin{verbatim}
P_i--p_{ij}--A_j,    Q_i--q_{ij}--B_j.
\end{verbatim}

Create row rails

\begin{verbatim}
P_i--x^L_{i,1}--y^L_{i,1}--x^L_{i,2}--...--y^L_{i,M}--Q_i
\end{verbatim}

and column rails

\begin{verbatim}
A_j--x^R_{j,1}--y^R_{j,1}--x^R_{j,2}--...--y^R_{j,M}--B_j.
\end{verbatim}

For every occurrence e=(i,j,s), create vertices $a_e$,$b_e$,$c_e$,$d_e$,$z_e$. Add paths

\begin{verbatim}
x^L_{i,s}--a_e--c_e--x^R_{j,s},
y^L_{i,s}--b_e--d_e--y^R_{j,s},
\end{verbatim}

and triangle $a_e$--$b_e$--$z_e$--$a_e$.

Attach sparse capacity-one selectors to each of the sets \{$p_{ij}$:j in[K]\} and \{$q_{ij}$:j in[K]\}, for each i. For each stage s, attach sparse capacity-one selectors to \{$c_e$:e in $E_s$\} and \{$d_e$:e in $E_s$\}, and a one-private-vertex capacity-(|$E_s$|-1) gadget to \{$z_e$:e in $E_s$\}. These port sets are pairwise disjoint.

Let $g=2K(K-1)+\sum_s(2|S_s|-1)$ be the total number of private vertices in the capacity gadgets. The threshold is $r_0=g+K+\sum_s|S_s|$.

\subsection{Completeness}

Suppose a permutation f hits every $S_s$. For each s choose one occurrence $e_s$=($i_s$,f($i_s$),s). Pack the triangles $a_e$,$b_e$,$z_e$ for all occurrences other than $e_s$. They contribute $sum_s$(|$S_s$|-1) cycles.

For a fixed row i and its paired column f(i), use the selector paths through $p_{i,f(i)}$ and $q_{i,f(i)}$, the two rails, and the cross paths of its selected occurrences $e_s$. If this row is selected at stages $s_1$<...<$s_h$, the start selector and first x cross-path close one cycle; consecutive y and next x cross-paths close h-1 intermediate cycles; the last y cross-path and end selector close a final cycle. When h=0, the start/end selectors and entire rails form one cycle. Thus each row contributes h+1 cycles. Summed over rows this is K+M cycles. They are vertex-disjoint because f is a permutation and the x and y endpoints of a stage are distinct.

Together these host cycles contribute K+$sum_s$|$S_s$| cycles. They use exactly one p and q selector port per row, exactly one c and d port per stage, and all but one z port per stage. By the matching property extend them by all g gadget 2-cycles, reaching $r_0$.

\subsection{Soundness}

Normalize a packing of at least $r_0$ cycles. Remove its g gadget 2-cycles, leaving a capacity-respecting host packing C of at least K+$sum_s$|$S_s$| cycles.

Let Z consist of all $P_i$,$Q_i$,$a_e$,$b_e$. Its size is 2K+2sum\_s|$S_s$|. The graph H-Z is a forest: its components are the interiors of the K row rails, the K column rails with selector ports and c/d leaves attached, and the isolated z vertices. Each vertex of Z has at most one neighbor in each component of H-Z. Thus a host cycle contains at least two vertices of Z. The number of cycles in C reaches |Z|/2, so equality holds: every host cycle contains exactly two Z vertices, and every Z vertex is covered.

The cycle through $P_i$ must use the initial row-rail edge and exactly one selector port $p_{i,f(i)}$. Using two p ports is forbidden by capacity one. At $A_{f(i)}$ this cycle must continue onto that column's rail: continuing through another p port would introduce another P vertex whose row component differs, requiring a third Z vertex before the cycle could close. Distinct P cycles use distinct A vertices, so f is a permutation. At the cut just before the first stage, each of these K distinct cycles pairs row i with column f(i), and each already contains one Z vertex before the cut.

More explicitly, the invariant is that on the past side of the cut there are K vertex-disjoint paths, path i joining row i to column f(i), and each such path contains exactly one Z vertex already encountered. Initially these paths contain $P_i$. Inductively assume this past-path pairing property at the cut before stage s. All stage-entry vertices $x^L_{i,s}$,$x^R_{j,s}$ belong to their incoming frontier cycles. Capacity on the z ports leaves some $z_e$ uncovered, for e=(i,j,s). Since $a_e$ is covered, it cannot use $z_e$. It also cannot share its cycle with $b_e$: joining through $b_e$ would add two new Z vertices to a cycle that already contains one Z vertex before this stage, because its other exit goes through a stage-entry vertex. Consequently $a_e$ uses the x cross-path through $c_e$, joining $x^L_{i,s}$ and $x^R_{j,s}$. If the two entry vertices belonged to different past paths, the cycle would contain $a_e$ and the two distinct earlier Z vertices, contradicting the exactly-two-Z bound. Thus they belong to the same past path and j=f(i). Therefore $S_s$ is hit.

The used $c_e$ saturates the stage's capacity-one c gadget. For any other occurrence e', $c_e$' is unavailable to host cycles. If $a_e$' were connected to its row entry, it would belong to an incoming frontier cycle, but the absence of $c_e$' would then force an additional new Z vertex $b_e$', contradicting the exactly-two-Z property. Hence $a_e$',$b_e$',$z_e$' form their local triangle. In particular every other z port is used, and e is the unique occurrence whose triangle is broken. Also $b_e$ cannot use $a_e$ or $z_e$, so it must use the y cross-path through $d_e$. Its row/column exits give the next past path pairing row i with column f(i), whose unique earlier Z vertex is now $b_e$; the x crossing has closed the old path's cycle. All other rails pass straight through the stage, because their attached a/b vertices are already occupied in triangles. Thus the same permutation pairing holds at the next cut. This establishes the induction and proves that f hits every set.

The last-stage argument proves its hit without requiring a further stage; its y cross-path reaches the final selector region. No uniqueness of a cell in a row of $S_s$ was used: capacity one is across all c ports of the stage and all d ports of the stage, independently of row multiplicity.

\subsection{Explicit path decomposition and simple graphs}

Keep the  $4K$ global vertices $P_i,Q_i,A_j,B_j$ throughout. Process the two endpoint selector systems using one row at a time. One may retain all K ports of the current row together with its K-1 private vertices; this adds fewer than 2K vertices. Their edges to globals and all doubled edges are covered.

For stage s, retain its $4K$ rail vertices $x^L_{i,s},y^L_{i,s},x^R_{j,s},y^R_{j,s}$. During a transition retain the preceding and next stage rail vertices together, covering all interstage rail edges; this contributes at most $8K$ in addition to the globals. Within the stage, order all occurrences in any fixed order. Process their five private host vertices a,b,c,d,z one occurrence at a time. Retain the previous occurrence's c,d ports until the next occurrence has been introduced, and introduce the two sparse-selector private vertices joining the previous and current c/d ports. Keep the single z-gadget center during the whole stage. A bound of 20 additional vertices beyond the global/rail boundary suffices. Forget the previous c,d and their selector-private vertices after these edges are covered; retain the current c,d for the next occurrence. Every occurrence interval is consecutive.

Thus supplied width is at most $12K+20$, independent of |$S_s$| and M. The number of vertices, edges and bags is polynomial in K+M+$sum_s$|$S_s$|.

Finally replace every edge, including each parallel copy, by a length-three path with two fresh internal vertices. Every resulting cycle corresponds to a multigraph cycle and vice versa; internally disjoint subdivisions preserve vertex-disjoint packing size. A supplied path decomposition extends by processing the two internal vertices at a bag containing the original endpoints, adding at most two to width. This yields a simple undirected graph with width at most $12K+22$. The start and end selector paths join all row and column rails. Each remaining vertex attaches to a rail or a selector, so the graph is connected.


\section{Reducing the degree}




Start with a simple graph G and a supplied path decomposition of width w. Order original edges by chosen covering bags, breaking ties consistently. At each nonisolated vertex v of degree d, build a backbone path with d vertices, numbered by this incident-edge order. At position i attach a port $p_i$ to its backbone vertex and a companion $x_i$ to $p_i$. Route each original edge through the two corresponding ports. For i from 1 through d-2, make a private center $u_i$ adjacent to $p_j$ and $x_j$ for j in \{i,i+1,i+2\}. Ignore isolated vertices. The ports have degree at most 2+1+3=6, companions degree at most 4, and centers degree six. All graphs at this stage are simple. For d\(\le\)2 there are no centers.

For d\(\ge\)2, any d-2 ports can be matched to the d-2 centers since their ordered indices $j_i$ satisfy i\(\le\)$j_i$\(\le\)i+2. A match supplies the triangle $u_i$ $p_{j_i}$ $x_{j_i}$. Thus all centers can be covered by disjoint triangles provided at most two ports and their companions are blocked in the gadget. A packing of G lifts by traversing the backbone between its two ports for each used original vertex. At most two ports are blocked in each gadget, and no companions are used, so the lifted packing gains the offset $sum_v$ max(d(v)-2,0).

For the reverse, begin with a maximum packing and normalize one gadget at a time. Remove all c cycles meeting its centers. Let a count ports used by kept cycles. Each kept occupied port uses its backbone and external edge since companions are leaves once centers are absent. Each kept cycle touching this gadget crosses its cut at least twice. It uses an even number of ports, so a is even. If a\(\le\)2 we can add d-2 triangles, which pays for the c removed cycles because they used distinct centers.

Suppose a\(\ge\)4. Any removed cycle that used only one center in the current gadget must have used at least one port. In the graph omitting ports and all other centers, a lone center and the companions adjacent to it form a star with no route to the outside. Such cycles used distinct ports available outside the kept packing. Thus if $c_1$ of them used one center, then $c_1$\(\le\)d-a, and $c_1$+2(c-$c_1$)\(\le\)d-2. It follows that c\(\le\)d-1-a/2. If fewer than a/2 kept cycles touch the gadget, delete them all. Otherwise exactly a/2 touch it, each at two ports, and we can retain any one and delete the others. In both cases deleting at most a/2-1 extra cycles leaves at most two occupied ports. Neither kept cycles nor cycles in previously normalized gadgets occupy the remaining companions. Matching adds all d-2 triangles, thereby paying for the removals. The case d\(\le\)2 needs no normalization.

New triangles are wholly local, so no previously installed triangles are disrupted. After all normalizations, remove the offset triangles. The remaining packing has at most two occupied ports per gadget and no private centers in use. Each visited gadget is traversed once by one cycle along a path between its two occupied ports. Contract those traversals to recover disjoint cycles in G. Therefore the offset identity holds exactly.

\subsection{Path decomposition}

Between edge events retain the newest backbone vertex and the newest two port-companion pairs per active original vertex. At an event introduce the new backbone vertex, its port and companion, and the center for the latest complete triple at both endpoints. This covers the new external edge, the backbone edges, and all six adjacencies of either center. Then drop the old backbone, the oldest port-companion pair when there are three, and each newly introduced center. All lifetimes are consecutive because an original vertex's events occur within its original bag interval. There are at most $5(w+1)+8$ vertices in a bag. The output graph and supplied decomposition have polynomial length. The transformation preserves connectivity.


\section{Forcing a spanning bipartite cover}
We use the conversions of Cygan et al.~\cite[Lemmas C.21 and C.22]{ccpre}
to preserve the packed cycle count up to additive offsets.


Start from a connected simple undirected Cycle Packing input of maximum degree six. The following two conversions use loopless digraphs, with a pair of opposite arcs permitted and a directed two-cycle counted as a cycle. For an old edge \(uv\), introduce private vertices \(x_{uv},x_{vu},z\). Add the directed paths \texttt{u -> x\_uv -> v} and \texttt{v -> x\_vu -> u}. Join each transit vertex to \(z\) by both arcs. The directed packing optimum is the old optimum plus \(|E(G)|\). Any old packing can be oriented and routed using at most one transit vertex per edge, leaving the other to form a two-cycle with \(z\).

For the reverse, normalize an optimum directed packing one edge gadget at a time. If its \(z\) is not already in a private two-cycle, replace the cycle containing \(z\), when present, by a two-cycle with a transit vertex on that cycle. If \(z\) is free, choose a free transit vertex if possible. When both are occupied, delete one cycle meeting a transit vertex and use it in the new two-cycle. In every case the count stays the same or increases, and no previously installed private two-cycle is affected. After the private two-cycles are removed, at most one transit vertex per old edge remains in use. Every used transit vertex now relays an original edge and every remaining cycle visits distinct original vertices. A projected two-cycle would use both transit vertices for the same edge, so all projections are simple undirected cycles of length at least three. Place each edge's three new vertices in one copy of a bag containing both endpoints. Width increases by at most three.

For every vertex \(v\) of this digraph, add a private directed triangle \texttt{a\_v -> b\_v -> c\_v -> a\_v} and arcs \texttt{v -> a\_v} and \texttt{c\_v -> v}. In a spanning cover the arcs $a_v\to b_v$ and $b_v\to c_v$ are forced. The next arc either closes the triangle at $a_v$ or goes to $v$, which then must have its outgoing arc to $a_v$. Conversely, each directed packing extends by choosing the triangle when $v$ is already covered and the four-cycle otherwise. Increase the threshold by the order of the digraph before this step and increase width by at most three. The maximum in-degree and out-degree after this step are at most seven. Original vertices start with each degree at most six, and transit and center vertices start with each at most two. The wrapper adds at most one arc of each orientation to each existing vertex.

Split each directed vertex \(v\) into two terminals (v\_in,v\_out) joined by a path of three edges with two private degree-two internal vertices. For each directed arc \texttt{u -> v} add the edge \texttt{u\_out -- v\_in}. Color all incoming terminals on one side and all outgoing terminals on the other. The internal paths extend this coloring. Every spanning undirected cover uses all internal path edges, since their internal vertices have degree two. It therefore uses exactly one external edge at each terminal, corresponding to one incoming and one outgoing selected arc at each old vertex. This gives an exact correspondence of covers, preserving the number of components. Simplicity follows from the absence of duplicate arcs, and the underlying graph stays connected at each step. Its maximum degree is at most eight. Lift each bag to the two terminals for every vertex, and add the two internal vertices in a copied bag for each old vertex. The supplied width is at most \texttt{2(t+6)+3 = 2t+15} for input width t, with polynomial complete output length.


\section{Running time}

Each member of the family has final supplied width
$p=O(K)=O(m/\log m)$ and polynomial complete length in the encoded
formula length. Take a fixed integer $A$ such that $\phi(p)\le Am$
for all output members, increasing it for small cases.
For a hypothetical algorithm choose $j\ge 4Ak$ computably.
For $p\ge T(j)$ we have $h(p)\le m/(4k)$.
Since $h$ is finite and effectively sublinear in $\phi$, there is
a fixed integer $C$ with $h(u)\le \phi(u)+C$ for every $u$.
At $p<T(j)$ the exponent is at most $\phi(T(j))+C$.
This provides a computable additive bound in $k$ on the small
width contributions. The effective enumeration exponent is
also at most $m/(4k)$ plus a computable additive term in $k$,
by the explicit envelope from the grouping construction.
With $m=k\lceil\log_2 n\rceil$, the total time is a computable
factor in $k$ times a fixed polynomial in $n$. Evaluating the
family of cover queries thus decides Clique in deterministic FPT
time by clique size, a contradiction.



\begin{thebibliography}{9}
\bibitem{df}
R. G. Downey and M. R. Fellows.
Fixed-parameter tractability and completeness II: On completeness for W[1].
\emph{Theoretical Computer Science}, 141(1--2):109--131, 1995.
\href{https://doi.org/10.1016/0304-3975(94)00097-3}{doi:10.1016/0304-3975(94)00097-3}.

\bibitem{lms}
D. Lokshtanov, D. Marx, and S. Saurabh.
Slightly superexponential parameterized problems.
\emph{SIAM Journal on Computing}, 47(3):675--702, 2018.
\href{https://doi.org/10.1137/16M1104834}{doi:10.1137/16M1104834}.
\bibitem{cc}
M. Cygan, J. Nederlof, M. Pilipczuk, M. Pilipczuk, J. M. M. van Rooij,
and J. O. Wojtaszczyk.
Solving connectivity problems parameterized by treewidth in single
exponential time.
\emph{ACM Transactions on Algorithms}, 18(2), Article 17, 2022.
\href{https://doi.org/10.1145/3506707}{doi:10.1145/3506707}.

\bibitem{ccpre}
M. Cygan et al.
Solving connectivity problems parameterized by treewidth in single
exponential time. \href{https://arxiv.org/abs/1103.0534v1}{arXiv:1103.0534v1}.
Appendix C.2--C.4.
\end{thebibliography}
\end{document}
